\documentclass{article}
\usepackage{amsmath, amssymb, amsthm}
\usepackage{hyperref}
\usepackage{geometry}
\geometry{margin=1in}

\title{Erd\H{o}s Problem \#478}
\date{Last updated: 2025-08-31}
\author{Source: erdosproblems.com}

\begin{document}
\maketitle

\noindent\textbf{Status:} OPEN \quad \textbf{No prize}

\noindent\textbf{Tags:} number theory, factorials

\medskip

\noindent\textbf{Problem Statement:}

Let $p$ be a prime and\[A_p = \{ k! \pmod{p} : 1\leq k&#60;p\}.\]Is it true that\[\lvert A_p\rvert \sim (1-\tfrac{1}{e})p?\]




Since $A_p/A_p=\{1,\ldots,p-1\}$ it follows that $\lvert A_p\rvert \gg p^{1/2}$. The best known lower bound is due to Grebennikov, Sagdeev, Semchankau, and Vasilevskii \cite{GSSV24},\[\lvert A_p\rvert \geq (\sqrt{2}-o(1))p^{1/2},\]which follows from proving that $\lvert A_pA_p\rvert=(1+o(1))p$. Wilson's theorem implies $(p-2)!\equiv 1\pmod{p}$, and hence $\lvert A_p\rvert\leq p-2$. Results on $\lvert A_p\rvert$ on average were obtained by Klurman and Munsch \cite{KlMu17}. It is even open whether there are infinitely many primes $p$ such that $\lvert A_p\rvert=p-2$: in other words, such that $2!,\ldots,(p-1)!$ are all distinct modulo $p$ (such primes are sometimes called socialist primes). Rokowska and Schinzel \cite{RoSc60} have proved that any socialist prime is $5\pmod{8}$, and must satisfy $-(\frac{5}{p})=(\frac{-23}{p})=1$. Moreover, the congruence class missing from $2!,\ldots,(p-1)!$ must be $-(\frac{p-1}{2})!$. The only known socialist prime is $p=5$. Trudgian \cite{Tr13} has shown there are no other socialist primes $&#60;10^9$. This was extended to $10^{11}$ by Andreji\'{c} and Tatarevic \cite{AnTa16}. This is also discussed in problems A2 and F11 of Guy's collection \cite{Gu04}.




References

[AnTa16] V. Andreji\'{c} and M. Tatarevic, On distinct residues of factorials . arXiv:1603.04086 (2016).

[GSSV24] Grebennikov, Alexandr and Sagdeev, Arsenii and Semchankau, Aliaksei and Vasilevskii, Aliaksei, On the sequence {$n! \bmod p$} . Rev. Mat. Iberoam. (2024), 637--648.

[Gu04] Guy, Richard K., Unsolved problems in number theory . (2004), xviii+437.

[KlMu17] Klurman, Oleksiy and Munsch, Marc, Distribution of factorials modulo {$p$} . J. Th\'{e}or. Nombres Bordeaux (2017), 169--177.

[RoSc60] Rokowska, B. and Schinzel, A., Sur un probl\`eme de {M}. {E}rd\H os . Elem. Math. (1960), 84--85.

[Tr13] T. Trudgian, There are no socialist primes less than $10^9$ . arXiv:1310.6403 (2013).

\noindent\textbf{Related OEIS sequences:} A210184


\bigskip
\noindent\small{Source: \url{https://www.erdosproblems.com/478}}
\end{document}
